\cleardoublepage
\section*{Abstract}

\par Image focus detection aims to determine whether local image regions are
clear or blurred, and it is important for image quality analysis, autofocus
assistance, and subsequent vision tasks. This thesis studies patch-based image
focus detection under limited data conditions and builds a unified experimental
pipeline covering image preprocessing, grayscale conversion, $32\times32$ patch
division, automatic labeling, model training, and prediction evaluation. To
reduce the cost of manual patch-level annotation, this thesis designs a
rule-based automatic labeling method by combining three kinds of sharpness
features, namely Laplacian, Sobel, and fast Fourier transform (FFT), and
organizes both binary and three-class classification tasks based on these
labels. Using a dataset constructed from $350$ original images, this thesis
systematically compares the performance of decision trees, random forests,
support vector machines based on Nystr\"{o}m kernel approximation, and a
lightweight convolutional neural network (CNN) under the same data split and
unified evaluation metrics.

\par The experimental results show that, under the current task setting and data
organization of this thesis, increasing the training sample ratio from $20\%$ to
$100\%$ brings only limited performance improvement, while the training time
keeps increasing. This indicates that the marginal benefit of data expansion is
limited within the scope of this study. Specifically, the lightweight CNN
achieves the best overall result in the binary classification task, reaching an
F1-score of $0.84$ for the clear class and a macro F1 of $0.91$ with $100\%$
training samples. In the three-class classification task, the random forest
obtains the highest F1-score of $0.83$ for the clear class, while the SVM
maintains a relatively stable macro F1 of about $0.89$. Overall, traditional
machine learning methods, especially random forest and SVM, are able to learn
competitive decision boundaries under a unified grayscale patch representation
and provide relatively stable results with lower training cost. Based on these
findings, this thesis argues that for patch-level focus detection tasks with
relatively clear local boundaries and limited samples, deep learning is not the
only approach that should be considered first. When the performance gap is
limited, traditional methods can still achieve a good balance between
recognition performance and computational efficiency.
